% !TEX program = xelatex
\documentclass[12pt,a4paper]{article}
\usepackage{fontspec}
\usepackage{amsmath,amssymb,amsfonts}
\usepackage{graphicx}
\usepackage{float}
\usepackage{subcaption}
\usepackage{xcolor}
\usepackage{listings}
\usepackage{hyperref}
\usepackage{booktabs}
\usepackage{geometry}
\usepackage{titlesec}
\usepackage{algorithm}
\usepackage{algpseudocode}
\usepackage{algorithmicx}
\usepackage{breqn}
\usepackage{bm}

% Set up fonts
\setmainfont[Ligatures=TeX]{Times New Roman}
\setsansfont{Arial}
\setmonofont{Courier New}

% Geometry settings
\geometry{left=2.5cm,right=2.5cm,top=2.5cm,bottom=2.5cm}
\usepackage{ragged2e}

\setlength{\parindent}{0pt}
\setlength{\parskip}{1em}

\emergencystretch=2em
\justifying

\floatname{algorithm}{Algorithm}
\renewcommand{\listalgorithmname}{List of Algorithms}

\lstset{
	language=Matlab,
	basicstyle=\ttfamily\footnotesize,
	keywordstyle=\color{blue},
	commentstyle=\color{green!60!black},
	numbers=left,
	numberstyle=\tiny\color{gray},
	stepnumber=1,
	numbersep=5pt,
	backgroundcolor=\color{gray!5},
	showspaces=false,
	showstringspaces=false,
	frame=single,
	tabsize=2,
	breaklines=true,
	captionpos=b
}

\newcommand{\eng}[1]{#1}
% ---- Tensor / operator macros ----
\newcommand{\bF}{\mathbf{F}}
\newcommand{\bI}{\mathbf{I}}
\newcommand{\bC}{\mathbf{C}}
\newcommand{\bsig}{\boldsymbol{\sigma}}
\newcommand{\beps}{\boldsymbol{\varepsilon}}
\newcommand{\bxi}{\boldsymbol{\xi}}
\newcommand{\bu}{\mathbf{u}}
\newcommand{\bv}{\mathbf{v}}
\newcommand{\bb}{\mathbf{b}}
\newcommand{\bn}{\mathbf{n}}
\newcommand{\bT}{\mathbf{T}}
\newcommand{\bB}{\mathbf{B}}
\newcommand{\bN}{\mathbf{N}}
\newcommand{\bL}{\mathbf{L}}
\newcommand{\bD}{\mathbf{D}}
\newcommand{\bK}{\mathbf{K}}
\newcommand{\bR}{\mathbf{R}}
\newcommand{\bx}{\mathbf{x}}
\newcommand{\Cz}{\mathbf{C}_0}
\DeclareMathOperator{\tr}{tr}
\DeclareMathOperator{\Div}{div}
\newcommand{\dd}{\,\mathrm{d}}
\newcommand{\Dt}{\Delta t}

\titleformat{\chapter}[display]
{\normalfont\huge\bfseries}{\chaptertitlename\ \thechapter}{20pt}{\Huge}

\algblockdefx[PHASE]{Phase}{EndPhase}[1]{\textbf{#1}}{\textbf{end}}

\begin{document}
	\setlength{\parindent}{0pt}

	% --- Title page ---
	\begin{titlepage}
		\centering
		{\LARGE\bfseries AN EXPLICIT USL MATERIAL POINT METHOD \\[.3cm]
		IMPLEMENTATION OF \\[.5cm] VARIATIONAL PHASE--FIELD FRACTURE}\\[3cm]

		{\large Technical Report}\\[1cm]

		{\large APOSTOLOPOULOS ORPHEAS}\\[0.5cm]
		{\large AM: 22030}\\[3cm]

		{\large SCHOOL OF MECHANICAL ENGINEERING}\\
		{\large NATIONAL TECHNICAL UNIVERSITY OF ATHENS}\\[5cm]

		{\large ATHENS, JUNE 2026}
	\end{titlepage}

	\tableofcontents
	\newpage

	% =====================================================================
	\section{Introduction and Scope}

	This report develops an \emph{explicit} solution scheme for variational
	phase--field fracture, formulated within the \emph{Update--Stress--Last}
	(\eng{USL}) Material Point Method (\eng{MPM}). The continuum phase--field
	theory is adopted essentially unchanged from the variational fracture
	framework summarised below, and the same symbols are retained throughout:
	the displacement field $\bu$, the scalar phase field (damage variable)
	$\phi\in[0,1]$, the degradation function $g(\phi)$, the crack surface density
	function $\gamma(\phi,\nabla\phi)$, the fracture toughness $G_c$, the length
	scale $\ell$, the normalisation constant $c_w$, the local crack geometry
	function $w(\phi)$, the (undamaged) elastic strain energy density $\psi^e$, and
	the history field $\mathcal{H}$.

	The original framework is intrinsically \emph{quasi--static}: the scalar
	micro--stress $\omega$ that is work conjugate to $\phi$ is purely
	\emph{energetic}, so the phase--field balance~$\nabla\cdot\bxi-\omega=0$ is an
	\emph{elliptic} equation that must be solved by an implicit (monolithic or
	staggered) Newton iteration at every time step. This is incompatible with an
	explicit MPM time integrator, in which all unknowns are marched forward in
	time without solving a linear system.

	\textbf{The central modification of this report} is the introduction of a
	\emph{dissipative} contribution to the phase--field micro--stress,
	\begin{equation}\label{eq:wdiss_intro}
		\boxed{\;\omega_{\mathrm{diss}} \;=\; \eta\,\dot{\phi} \;=\; \eta\,\frac{\dd\phi}{\dd t}\;,\qquad \eta\ge 0,\;}
	\end{equation}
	where $\eta$ is a (artificial) viscosity with units of
	$\mathrm{Pa\cdot s}$. As shown by Kakouris \& Triantafyllou and by Cheon \&
	Kim, this single term turns the damage sub--problem from an \emph{elliptic}
	into a \emph{parabolic} (Allen--Cahn type) equation, so the phase field can be
	advanced in time by an explicit forward--Euler update on a lumped pseudo--mass,
	exactly in the spirit of the explicit momentum update of the displacement
	problem. The result is a fully explicit, monolithic--in--time
	\eng{USL--MPM} algorithm in which displacement and phase field are updated
	within the \emph{same} time loop.

	The report is organised as follows. Section~\ref{sec:reg} recalls the
	phase--field regularisation of Griffith fracture. Section~\ref{sec:vw} restates
	the principle of virtual work, now augmented with the dissipative
	micro--force~\eqref{eq:wdiss_intro}, and derives the strong forms.
	Section~\ref{sec:const} gives the constitutive theory and the AT1/AT2 crack
	density choices. Section~\ref{sec:strong} collects the coupled strong forms,
	emphasising the parabolic nature of the regularised evolution law.
	Section~\ref{sec:disc} treats damage irreversibility, the weak form and the
	MPM space discretisation, and derives the \emph{explicit nodal phase--field
	update}. Finally, Section~\ref{sec:algo} embeds this update into the explicit
	\eng{USL--MPM} cycle and states the complete solution algorithm, followed by
	stability and time--step considerations.

	% =====================================================================
	\section{Phase--field Regularisation of Griffith Fracture}\label{sec:reg}

	Consider a solid of volume $V$ occupying a domain $\Omega\subset\mathbb{R}^n$
	($n\in\{1,2,3\}$) with external boundary $\partial\Omega$ and outward unit
	normal $\bn$. The elastic strain energy density of the undamaged solid is
	$\psi^e(\beps)$, a function of the (small) strain tensor $\beps$.

	Griffith's energy balance states that the variation of the total energy $\Pi$
	with respect to an incremental increase of the crack area $\dd A$ vanishes,
	\begin{equation}\label{eq:griffith_dPi}
		\frac{\dd\Pi}{\dd A} \;=\; \frac{\dd\psi^e(\beps)}{\dd A} \;+\; \frac{\dd W_c}{\dd A} \;=\; 0,
	\end{equation}
	where $W_c$ is the work required to create new surfaces and
	$G_c=\dd W_c/\dd A$ is the critical energy release rate. In variational form,
	\begin{equation}\label{eq:griffith_var}
		\Pi \;=\; \int_{\Omega}\psi^e(\beps)\dd V \;+\; \int_{\Gamma}G_c\dd\Gamma,
	\end{equation}
	with $\Gamma$ the (unknown) crack surface. The difficulty posed by the unknown
	$\Gamma$ is removed by introducing an auxiliary phase field $\phi$ that smears
	the sharp interface into a diffuse region: $\phi=0$ in intact material and
	$\phi=1$ in fully cracked material. The Griffith functional is then approximated
	by the regularised functional
	\begin{equation}\label{eq:Pi_ell}
		\Pi_\ell \;=\; \int_{\Omega}\bigl[\,g(\phi)\,\psi^e_0(\beps) \;+\; G_c\,\gamma(\phi,\ell)\,\bigr]\dd V,
	\end{equation}
	where $\ell$ is a length scale governing the size of the fracture process zone,
	$\psi^e_0$ is the elastic energy density of the undamaged solid, and
	$\gamma(\phi,\ell)$ is the crack surface density function. The degradation
	function $g(\phi)$ degrades the stiffness as $\phi\to 1$. As $\ell\to 0^+$ the
	functional $\Pi_\ell$ $\Gamma$--converges to the sharp Griffith
	functional~\eqref{eq:griffith_var}; for $\ell>0^+$ a finite material strength
	is introduced, so that $\ell$ becomes a material property governing strength,
	e.g.\ for plane stress
	\begin{equation}\label{eq:sigc}
		\sigma_c \;\propto\; \sqrt{\frac{G_c E}{\ell}} \;=\; \frac{K_{Ic}}{\sqrt{\ell}},
	\end{equation}
	with $K_{Ic}$ the fracture toughness and $E$ Young's modulus.

	% =====================================================================
	\section{Principle of Virtual Work and Dissipative Micro--force}\label{sec:vw}

	The primal kinematic variables are the displacement field $\bu$ and the phase
	field $\phi$. Restricting attention to small strains and isothermal conditions,
	the strain tensor is
	\begin{equation}\label{eq:strain}
		\beps \;=\; \tfrac{1}{2}\bigl(\nabla\bu^{T}+\nabla\bu\bigr).
	\end{equation}
	Let $\delta$ denote virtual quantities. The Cauchy stress $\bsig$ is work
	conjugate to $\beps$, the traction $\bT$ is work conjugate to $\bu$, the scalar
	stress--like quantity $\omega$ is work conjugate to $\phi$, and the phase--field
	micro--stress vector $\bxi$ is work conjugate to $\nabla\phi$. Because the phase
	field is driven by the displacement problem alone, no external traction is
	associated with $\phi$. In the absence of body forces the principle of virtual
	work reads
	\begin{equation}\label{eq:vw}
		\int_{\Omega}\bigl\{\bsig:\delta\beps \;+\; \omega\,\delta\phi \;+\; \bxi\cdot\nabla\delta\phi\bigr\}\dd V
		\;=\; \int_{\partial\Omega}\bigl(\bT\cdot\delta\bu\bigr)\dd S.
	\end{equation}

	\paragraph{Energetic and dissipative parts of the micro--stress.}
	The key step for an explicit scheme is to allow the phase--field micro--stress
	$\omega$ to possess \emph{both} an energetic and a dissipative part,
	\begin{equation}\label{eq:w_split}
		\omega \;=\; \omega_{\mathrm{en}} \;+\; \omega_{\mathrm{diss}},
		\qquad
		\omega_{\mathrm{en}} \;=\; \frac{\partial\psi}{\partial\phi},
		\qquad
		\omega_{\mathrm{diss}} \;=\; \eta\,\dot{\phi},
	\end{equation}
	where $\psi$ is the free energy density (defined in
	Section~\ref{sec:const}) and $\eta\ge 0$ is the phase--field viscosity. The
	dissipative part $\omega_{\mathrm{diss}}=\eta\dot{\phi}$ is the minimal
	rate--dependent contribution that satisfies the dissipation inequality
	$\omega_{\mathrm{diss}}\,\dot{\phi}=\eta\,\dot{\phi}^2\ge 0$ and vanishes in the
	rate--independent (quasi--static) limit $\eta\to 0$. The micro--stress
	vector $\bxi$ is taken purely energetic,
	$\bxi=\partial\psi/\partial\nabla\phi$.

	\paragraph{Local balance laws (strong form).}
	Since~\eqref{eq:vw} must hold for an arbitrary domain $\Omega$ and for any
	kinematically admissible variation, the fundamental lemma of the calculus of
	variations yields the local force balances
	\begin{equation}\label{eq:balances}
		\nabla\cdot\bsig \;=\; \mathbf{0},
		\qquad
		\nabla\cdot\bxi \;-\; \omega \;=\; 0
		\qquad\text{in }\Omega,
	\end{equation}
	with natural boundary conditions
	\begin{equation}\label{eq:bcs}
		\bsig\cdot\bn \;=\; \bT,
		\qquad
		\bxi\cdot\bn \;=\; 0
		\qquad\text{on }\partial\Omega.
	\end{equation}
	Substituting the split~\eqref{eq:w_split} into the phase--field balance
	$\nabla\cdot\bxi-\omega=0$ gives the \emph{rate} form
	\begin{equation}\label{eq:rate_balance}
		\boxed{\;\eta\,\dot{\phi} \;=\; \nabla\cdot\bxi \;-\; \omega_{\mathrm{en}}\;.}
	\end{equation}
	Equation~\eqref{eq:rate_balance} is the cornerstone of the explicit scheme:
	for $\eta>0$ it is a first--order--in--time evolution law that can be
	integrated by forward Euler, whereas for $\eta=0$ it collapses to the original
	elliptic stationarity condition $\nabla\cdot\bxi-\omega_{\mathrm{en}}=0$.
	The dynamic momentum balance used by the explicit MPM replaces the first
	of~\eqref{eq:balances} by $\rho\,\dot{\bv}=\nabla\cdot\bsig+\rho\,\bb$, with
	$\rho$ the density and $\bb$ the body force per unit mass.

	% =====================================================================
	\section{Constitutive Theory}\label{sec:const}

	The free energy density of the solid is additively split into an elastic
	(bulk) part and a fracture part,
	\begin{equation}\label{eq:psi}
		\psi(\beps,\phi,\nabla\phi) \;=\; \psi^e \;+\; \psi^f
		\;=\; g(\phi)\,\tfrac{1}{2}\,\beps:\Cz:\beps \;+\; G_c\,\gamma(\phi,\nabla\phi),
	\end{equation}
	where $\Cz$ is the linear elastic stiffness tensor and
	$\psi^e_0=\tfrac{1}{2}\beps:\Cz:\beps$ is the undamaged elastic energy
	density. The Cauchy stress follows as
	\begin{equation}\label{eq:cauchy}
		\bsig \;=\; \frac{\partial\psi}{\partial\beps} \;=\; g(\phi)\,(\Cz:\beps) \;=\; g(\phi)\,\bsig_0,
		\qquad \bsig_0 \equiv \Cz:\beps ,
	\end{equation}
	with $\bsig_0$ the effective (undamaged) stress.

	\paragraph{Degradation function.}
	The degradation function is continuous and monotonic with $g(0)=1$ and
	$g(1)=0$; the standard quadratic form is adopted,
	\begin{equation}\label{eq:g}
		g(\phi) \;=\; (1-\phi)^2,
		\qquad g'(\phi)=-2(1-\phi),
		\qquad g''(\phi)=2.
	\end{equation}

	\paragraph{Crack surface density (AT1 / AT2).}
	Following the Ambrosio--Tortorelli family, the crack surface density function
	and the crack surface $A$ are
	\begin{equation}\label{eq:gamma}
		A \;=\; \int_{\Omega}\gamma(\phi)\dd V
		\;=\; \int_{\Omega}\frac{1}{4c_w\ell}\Bigl(w(\phi) \;+\; \ell^2|\nabla\phi|^2\Bigr)\dd V,
	\end{equation}
	where the local crack geometry function $w(\phi)$ satisfies $w(0)=0$, $w(1)=1$,
	and the normalisation constant is
	\begin{equation}\label{eq:cw}
		c_w \;=\; \int_0^1\sqrt{w(\varphi)}\,\dd\varphi .
	\end{equation}
	Two standard choices are retained throughout this report:
	\begin{align}
		\textbf{AT2:}\quad & w(\phi)=\phi^2,\quad c_w=\tfrac12, &
		w'(\phi)&=2\phi, & w''(\phi)&=2, \label{eq:AT2}\\[2pt]
		\textbf{AT1:}\quad & w(\phi)=\phi,\quad\; c_w=\tfrac23, &
		w'(\phi)&=1, & w''(\phi)&=0. \label{eq:AT1}
	\end{align}
	The AT2 model degrades from the onset of loading, whereas the AT1 model
	exhibits a purely elastic regime prior to damage onset. The critical
	(homogeneous, $\nabla\phi=0$) failure stresses are
	\begin{equation}\label{eq:sigc_models}
		\sigma_c^{\,\mathrm{AT1}} \;=\; \sqrt{\frac{3EG_c}{8\ell}},
		\qquad
		\sigma_c^{\,\mathrm{AT2}} \;=\; \frac{3}{16}\sqrt{\frac{3EG_c}{\ell}} .
	\end{equation}

	\paragraph{Micro--stresses.}
	From~\eqref{eq:psi}, the energetic scalar micro--stress and the
	micro--stress vector are
	\begin{align}
		\omega_{\mathrm{en}} \;=\; \frac{\partial\psi}{\partial\phi}
		&\;=\; g'(\phi)\,\psi^e \;+\; \frac{1}{4c_w\ell}\,G_c\,w'(\phi), \label{eq:wen}\\[4pt]
		\bxi \;=\; \frac{\partial\psi}{\partial\nabla\phi}
		&\;=\; \frac{\ell}{2c_w}\,G_c\,\nabla\phi. \label{eq:xi}
	\end{align}
	Together with the dissipative part~\eqref{eq:w_split}, the total
	micro--stress is $\omega=g'(\phi)\psi^e+\tfrac{1}{4c_w\ell}G_c w'(\phi)+\eta\dot\phi$.

	% =====================================================================
	\section{Coupled Strong Forms}\label{sec:strong}

	Collecting~\eqref{eq:rate_balance}, \eqref{eq:wen} and~\eqref{eq:xi}, the
	regularised phase--field evolution law becomes the parabolic (Allen--Cahn type)
	equation
	\begin{equation}\label{eq:evolution}
		\boxed{\;
		\eta\,\dot{\phi}
		\;=\; \frac{\ell}{2c_w}G_c\,\nabla^2\phi
		\;-\; \frac{G_c}{4c_w\ell}\,w'(\phi)
		\;-\; g'(\phi)\,\psi^e(\beps)\;,}
	\end{equation}
	obtained from $\nabla\cdot\bxi=\tfrac{\ell}{2c_w}G_c\nabla^2\phi$. Equivalently,
	using $g'(\phi)=-2(1-\phi)$, the right--hand side expresses the competition
	between the stored elastic energy (which drives $\phi\to 1$) and the fracture
	energy (which penalises crack surface). In the quasi--static limit $\eta\to 0$,
	\eqref{eq:evolution} reduces to the stationarity condition of the original
	theory,
	\begin{equation}\label{eq:stationary}
		\frac{G_c}{2c_w\ell}\Bigl(\frac{w'(\phi)}{2}-\ell^2\nabla^2\phi\Bigr)
		\;+\; g'(\phi)\,\psi^e(\beps) \;=\; 0 .
	\end{equation}

	The complete coupled boundary/initial--value problem to be solved is therefore
	\begin{subequations}\label{eq:strongall}
	\begin{align}
		\rho\,\dot{\bv} &\;=\; \nabla\cdot\bigl[g(\phi)\,\bsig_0\bigr] \;+\; \rho\,\bb, & &\text{(momentum)} \label{eq:strong_mom}\\[2pt]
		\eta\,\dot{\phi} &\;=\; \frac{\ell}{2c_w}G_c\,\nabla^2\phi - \frac{G_c}{4c_w\ell}w'(\phi) - g'(\phi)\,\mathcal{H}, & &\text{(phase field)} \label{eq:strong_pf}\\[2pt]
		\bsig\cdot\bn &\;=\; \bT, \qquad \bxi\cdot\bn=\nabla\phi\cdot\bn=0, & &\text{(boundary)} \label{eq:strong_bc}
	\end{align}
	\end{subequations}
	where, anticipating Section~\ref{sec:disc}, the elastic energy
	$\psi^e$ driving the phase field has been replaced by the irreversible history
	field $\mathcal{H}$. The two equations are coupled through $g(\phi)$ (the phase
	field degrades the stress in the momentum balance) and through
	$\psi^e$/$\mathcal{H}$ (the elastic energy drives the phase field).

	% =====================================================================
	\section{Discretisation and Explicit Phase--field Update}\label{sec:disc}

	\subsection{Damage irreversibility and history field}

	To prevent crack healing, the phase field must be monotonically
	non--decreasing,
	\begin{equation}\label{eq:irrev}
		\phi_{t+\Dt} \;\ge\; \phi_t .
	\end{equation}
	Following Miehe \emph{et al.}, this is enforced by replacing the instantaneous
	elastic energy $\psi^e_0$ with a history field $\mathcal{H}$ that records the
	maximum elastic energy density attained over the loading history,
	\begin{equation}\label{eq:history}
		\mathcal{H} \;=\; \max_{\tau\in[0,t]}\psi^e_0(\tau),
		\qquad
		\psi^e_0-\mathcal{H}\le 0,\quad \dot{\mathcal{H}}\ge 0,\quad \dot{\mathcal{H}}\,(\psi^e_0-\mathcal{H})=0 .
	\end{equation}
	The crack driving force in the phase--field equation is thus $\mathcal{H}$, which
	is updated particle--wise (Section~\ref{sec:algo}) and is monotone by
	construction.

	\subsection{Weak form of the damage sub--problem}

	Adopting a total--Lagrangian description for the damage equation and using
	the dissipative micro--force~\eqref{eq:w_split}, the weak form of the
	phase--field balance~\eqref{eq:rate_balance} over the reference domain
	$\Omega_0$ reads, for any admissible test function $\delta\phi$,
	\begin{equation}\label{eq:weak_pf}
		\int_{\Omega_0}\Bigl\{\,g'(\phi)\,\mathcal{H}\,\delta\phi
		\;+\; \frac{G_c}{2c_w\ell}\Bigl[\tfrac{1}{2}w'(\phi)\,\delta\phi
		\;+\; \ell^2\,\nabla_0\phi\cdot\nabla_0\delta\phi\Bigr]
		\;+\; \eta\,\dot{\phi}\,\delta\phi\,\Bigr\}\dd V \;=\; 0 .
	\end{equation}
	The first three terms are exactly the energetic residual of the original
	quasi--static formulation; the last term, $\eta\dot\phi\,\delta\phi$, is the new
	dissipative contribution that renders the problem parabolic.

	\subsection{Spatial discretisation}

	Both the displacement and the phase field are interpolated with the same shape
	functions. In the MPM the shape functions are the \emph{background grid}
	functions $N_I$ (equivalently $\phi_I$ in some MPM notations); we write $N_I$ to
	avoid clashing with the phase field $\phi$. Nodal quantities are interpolated as
	\begin{equation}\label{eq:interp}
		\bu \;=\; \sum_{I} N_I\,\hat{\bu}_I,
		\qquad
		\phi \;=\; \sum_{I} N_I\,\hat{\phi}_I,
		\qquad
		\beps \;=\; \sum_{I}\bB_I^{u}\,\hat{\bu}_I,
		\qquad
		\nabla\phi \;=\; \sum_{I}\bB_I\,\hat{\phi}_I,
	\end{equation}
	with $\bB_I^{u}$ the strain--displacement matrix and $\bB_I=\nabla N_I$ the
	gradient of the nodal shape function. Inserting~\eqref{eq:interp}
	into the weak form~\eqref{eq:weak_pf} yields, at node $I$, the residual
	\begin{equation}\label{eq:res_pf}
		R_I^{\phi}
		\;=\; \int_{\Omega_0}\Bigl\{\,g'(\phi)\,N_I\,\mathcal{H}
		\;+\; \frac{G_c}{2c_w\ell}\Bigl[\tfrac{1}{2}w'(\phi)\,N_I
		\;+\; \ell^2\,\bB_I^{T}\,\nabla_0\phi\Bigr]
		\;+\; \eta\,\dot{\phi}\,N_I\,\Bigr\}\dd V \;=\; 0 .
	\end{equation}
	For reference, the (small--strain) displacement residual and the consistent
	tangents of the implicit formulation are
	\begin{align}
		\bR_I^{u} &\;=\; \int_{\Omega_0}\bigl[g(\phi)+\kappa\bigr]\,(\bB_I^{u})^{T}\bsig_0\dd V, \label{eq:res_u}\\
		\bK_{IJ}^{u} &\;=\; \int_{\Omega_0}\bigl[g(\phi)+\kappa\bigr](\bB_I^{u})^{T}\Cz\,\bB_J^{u}\dd V, \label{eq:Kuu}\\
		K_{IJ}^{\phi} &\;=\; \int_{\Omega_0}\Bigl\{\Bigl(g''(\phi)\,\mathcal{H}+\frac{G_c}{4c_w\ell}w''(\phi)\Bigr)N_I N_J + \frac{G_c\,\ell}{2c_w}\bB_I^{T}\bB_J\Bigr\}\dd V, \label{eq:Kpp}
	\end{align}
	where $\kappa\approx 10^{-7}$ is a small residual stiffness that prevents
	ill--conditioning when $\phi=1$. The explicit scheme below does \emph{not}
	assemble or invert $K^{\phi}$; it is reproduced only to make the connection with
	the implicit formulation explicit.

	\subsection{Explicit nodal phase--field update}\label{sec:explicit_update}

	Discretising the dissipative term in~\eqref{eq:res_pf} with a forward (explicit)
	Euler approximation of the rate,
	\begin{equation}\label{eq:fwd_euler}
		\dot{\phi} \;\approx\; \frac{\phi_I^{\,t+\Dt}-\phi_I^{\,t}}{\Dt},
	\end{equation}
	and lumping(eg row-sum rule) the viscous term onto a diagonal \emph{pseudo--mass}
	\begin{equation}\label{eq:pseudomass}
		m_I \;:=\; \int_{\Omega_0} N_I\dd V ,
	\end{equation}
	the nodal balance $R_I^{\phi}=0$ can be solved \emph{explicitly} for the
	updated nodal phase field. Writing the energetic part of the residual as
	$F_I$,
	\begin{equation}\label{eq:Fnode}
		F_I \;=\; \int_{\Omega_0}\Bigl\{\,g'(\phi^t)\,N_I\,\mathcal{H}
		\;+\; \frac{G_c}{2c_w\ell}\Bigl[\tfrac{1}{2}w'(\phi^t)\,N_I
		\;+\; \ell^2\,\bB_I^{T}\,\nabla_0\phi^{\,t}\Bigr]\Bigr\}\dd V,
	\end{equation}
	the update reads
	\begin{equation}\label{eq:nodal_update}
		\boxed{\;
		\phi_I^{\,t+\Dt} \;=\; \phi_I^{\,t} \;-\; \frac{\Dt}{\eta\,m_I}\,F_I \;.}
	\end{equation}

	\paragraph{MPM (particle) quadrature.}
	In the MPM the reference--domain integrals are evaluated using the material
	points as quadrature points, each carrying a reference volume $V_p^0$, so that
	$\int_{\Omega_0}(\,\cdot\,)\dd V\approx\sum_p(\,\cdot\,)_p\,V_p^0$ and the
	pseudo--mass becomes $m_I=\sum_p N_I(\bx_p)\,V_p^0$. The explicit nodal update
	in fully discrete MPM form is
	\begin{equation}\label{eq:mpm_update}
		\boxed{\;
		\phi_I^{\,t+\Dt} \;=\; \phi_I^{\,t}
		\;-\; \frac{\Dt}{\eta\,m_I}\sum_{p}\Bigl[\,g'(\phi_p^{t})\,N_I(\bx_p)\,\mathcal{H}_p
		\;+\; \frac{G_c}{2c_w\ell}\Bigl(\tfrac12 w'(\phi_p^{t})\,N_I(\bx_p)
		\;+\; \ell^2\,\nabla N_I(\bx_p)\cdot\nabla\phi^{\,t}\Bigr)\Bigr]V_p^0 \;.}
	\end{equation}
	Here $\phi_p^{t}=\sum_I N_I(\bx_p)\phi_I^{t}$ and
	$\nabla\phi^{t}=\sum_I\nabla N_I(\bx_p)\phi_I^{t}$ are the phase field and its
	gradient evaluated at the particle, and $\mathcal{H}_p$ is the particle history
	field. After the grid update~\eqref{eq:mpm_update}, the particle phase field is
	recovered by mapping back from the grid,
	\begin{equation}\label{eq:p2phi}
		\phi_p^{\,t+\Dt} \;=\; \sum_I N_I(\bx_p)\,\phi_I^{\,t+\Dt},
	\end{equation}
	and the irreversibility constraint~\eqref{eq:irrev} is enforced
	point--wise, $\phi_p^{t+\Dt}\leftarrow\max(\phi_p^{t},\phi_p^{t+\Dt})$.

	\paragraph{Remark (mapping with the original symbols).}
	The structure of~\eqref{eq:mpm_update} is identical to the explicit
	\eng{ULMPM--PFM} update of Kakouris \& Triantafyllou / Cheon \& Kim, under the
	symbol correspondence $d\leftrightarrow\phi$, $\omega(d)\leftrightarrow g(\phi)$,
	$\breve{Y}\leftrightarrow\mathcal{H}$, $\alpha(d)\leftrightarrow w(\phi)$,
	$b\leftrightarrow\ell$, $c_\alpha\leftrightarrow c_w$ (up to the chosen
	normalisation of $\gamma$), and the viscosity $\xi\leftrightarrow\eta$. We retain
	the symbols of the variational theory throughout.

	% =====================================================================
	\section{Total Solution Algorithm: Explicit \eng{USL--MPM} with Phase Field}\label{sec:algo}

	The phase--field update~\eqref{eq:mpm_update} is now embedded in the explicit
	\eng{USL} (Update--Stress--Last) MPM cycle. In USL, the particle stresses are
	updated \emph{last}, i.e.\ after the grid momenta have been advanced and mapped
	back to the particles. Each material point $p$ carries position $\bx_p$,
	velocity $\bv_p$, stress $\bsig_p$ (or deformation gradient $\bF_p$), volumes
	$V_p^0,V_p$, mass $m_p$, density $\rho_p$, phase field $\phi_p$ and history
	$\mathcal{H}_p$.

	The phase field is advanced within the same particle--to--grid / grid--to--particle
	machinery as the momentum problem:
	\begin{enumerate}
		\item In the \emph{particle--to--grid} (P2G) sweep, the pseudo--mass
		$m_I=\sum_p N_I(\bx_p)V_p^0$ and the energetic damage force $F_I$
		of~\eqref{eq:Fnode} are accumulated, alongside the mechanical nodal mass,
		momentum and forces.
		\item The grid phase field is advanced explicitly by~\eqref{eq:mpm_update},
		in lock--step with the explicit momentum update.
		\item In the \emph{grid--to--particle} (G2P) sweep, the particle phase field is
		updated by~\eqref{eq:p2phi}, the history field $\mathcal{H}_p$ is refreshed
		from the current elastic energy, and the particle stresses are updated last
		using the degraded constitutive law $\bsig_p=g(\phi_p)\bsig_{0,p}$.
	\end{enumerate}

	The complete procedure is given in Algorithm~\ref{alg:uslmpm_pf}. Lines that are
	new relative to the standard \eng{USL--MPM} (Algorithm~1 of the reference
	formulation) are those handling the phase field: the pseudo--mass and damage
	force accumulation in the P2G sweep, the explicit grid damage update, and the
	particle damage / history update in the G2P sweep.

	\begin{algorithm}[H]
		\caption{Solution procedure of explicit \eng{USL--MPM} with phase--field fracture}
		\label{alg:uslmpm_pf}
		\begin{algorithmic}[1]
			\State \textbf{Initialisation:} set up Cartesian grid; $t\leftarrow 0$
			\State Set up particle data: $\bx_p^0,\bv_p^0,\bsig_p^0,\bF_p^0,V_p^0,m_p,\rho_p^0,\phi_p^0,\mathcal{H}_p^0$
			\While{$t<t_f$}
				\State \textbf{Reset grid quantities:} $m_I\!\leftarrow\!0,\;(m\bv)_I\!\leftarrow\!\mathbf{0},\;\mathbf{f}_I^{\mathrm{ext}}\!\leftarrow\!\mathbf{0},\;\mathbf{f}_I^{\mathrm{int}}\!\leftarrow\!\mathbf{0},\;m_I^{\phi}\!\leftarrow\!0,\;F_I\!\leftarrow\!0$
				\Phase{Particle$\,\to\,$grid mapping (P2G)}
					\State Nodal mass: $m_I^{t}=\sum_p N_I(\bx_p^t)\,m_p$
					\State Nodal momentum: $(m\bv)_I^{t}=\sum_p N_I(\bx_p^t)\,(m\bv)_p$
					\State External force: $\mathbf{f}_I^{\mathrm{ext},t}=\sum_p N_I(\bx_p^t)\,m_p\,\bb(\bx_p)$
					\State Internal force: $\mathbf{f}_I^{\mathrm{int},t}=-\sum_p V_p^{t}\,\bsig_p^{t}\,\nabla N_I(\bx_p^t)$ \Comment{$\bsig_p^t=g(\phi_p^t)\bsig_{0,p}^t$ degraded}
					\State Nodal force: $\mathbf{f}_I^{t}=\mathbf{f}_I^{\mathrm{ext},t}+\mathbf{f}_I^{\mathrm{int},t}$
					\State Pseudo--mass: $m_I^{\phi}=\sum_p N_I(\bx_p^t)\,V_p^0$
					\State Damage force: $F_I=\sum_p\Bigl[g'(\phi_p^t)N_I(\bx_p^t)\mathcal{H}_p^t+\dfrac{G_c}{2c_w\ell}\bigl(\tfrac12 w'(\phi_p^t)N_I(\bx_p^t)+\ell^2\nabla N_I(\bx_p^t)\!\cdot\!\nabla\phi^t\bigr)\Bigr]V_p^0$
				\EndPhase
				\State \textbf{Update momenta:} $(m\bv)_I^{\,t+\Dt}=(m\bv)_I^{t}+\mathbf{f}_I^{t}\,\Dt$
				\State \textbf{Update grid phase field:} $\phi_I^{\,t+\Dt}=\phi_I^{t}-\dfrac{\Dt}{\eta\,m_I^{\phi}}\,F_I$
				\State \textbf{Fix Dirichlet nodes} $I$: e.g.\ $(m\bv)_I^{t}=\mathbf{0}$ and $(m\bv)_I^{\,t+\Dt}=\mathbf{0}$; prescribe $\phi_I$ on damage--BC nodes
				\Phase{Grid$\,\to\,$particle mapping (G2P)}
					\State Nodal velocities: $\bv_I^{t}=(m\bv)_I^{t}/m_I^{t}$, \;\; $\bv_I^{\,t+\Dt}=(m\bv)_I^{\,t+\Dt}/m_I^{t}$
					\State Update particle velocities (PIC/FLIP blend): \par
					\hskip1.2em $\bv_p^{\,t+\Dt}=\alpha\bigl(\bv_p^{t}+\sum_I N_I(\bx_p^t)[\bv_I^{\,t+\Dt}-\bv_I^{t}]\bigr)+(1-\alpha)\sum_I N_I(\bx_p^t)\,\bv_I^{\,t+\Dt}$
					\State Update particle positions: $\bx_p^{\,t+\Dt}=\bx_p^{t}+\Dt\sum_I N_I(\bx_p^t)\,\bv_I^{\,t+\Dt}$
					\State Update particle phase field: $\phi_p^{\,t+\Dt}=\sum_I N_I(\bx_p^t)\,\phi_I^{\,t+\Dt}$
					\State Enforce irreversibility: $\phi_p^{\,t+\Dt}\leftarrow\max\!\bigl(\phi_p^{t},\,\phi_p^{\,t+\Dt}\bigr)$
					\State Velocity gradient: $\bL_p^{\,t+\Dt}=\sum_I \nabla N_I(\bx_p^t)\,\bv_I^{\,t+\Dt}$
					\State Deformation gradient: $\bF_p^{\,t+\Dt}=(\bI+\bL_p^{\,t+\Dt}\Dt)\,\bF_p^{t}$
					\State Update volume: $V_p^{\,t+\Dt}=\det(\bF_p^{\,t+\Dt})\,V_p^0$
					\State Rate of deformation: $\bD_p^{\,t+\Dt}=\tfrac12\bigl(\bL_p^{\,t+\Dt}+(\bL_p^{\,t+\Dt})^{T}\bigr)$
					\State Strain increment: $\Delta\beps_p=\Dt\,\bD_p^{\,t+\Dt}$
					\State \textbf{Update stress (last):} effective $\bsig_{0,p}^{\,t+\Dt}=\bsig_{0,p}^{t}+\Cz:\Delta\beps_p$; degraded $\bsig_p^{\,t+\Dt}=g(\phi_p^{\,t+\Dt})\,\bsig_{0,p}^{\,t+\Dt}$
					\State Update elastic energy and history: $\psi^e_{0,p}=\tfrac12\beps_p:\Cz:\beps_p$, \;\; $\mathcal{H}_p^{\,t+\Dt}=\max\!\bigl(\mathcal{H}_p^{t},\,\psi^e_{0,p}\bigr)$
				\EndPhase
				\State \textbf{Advance time:} $t\leftarrow t+\Dt$
			\EndWhile
		\end{algorithmic}
	\end{algorithm}

	\subsection{Time--step and stability considerations}

	Two explicit stability restrictions act simultaneously, and the time step must
	respect the smaller of the two.

	\paragraph{Elastodynamic (CFL) limit.}
	The momentum update is conditionally stable and limited by the
	Courant--Friedrichs--Lewy condition,
	\begin{equation}\label{eq:cfl}
		\Dt \;\le\; \Dt_{\mathrm{crit}}^{\,u} \;=\; \frac{h}{c_d},
		\qquad
		c_d \;=\; \sqrt{\frac{\lambda+2\mu}{\rho}},
	\end{equation}
	with $h$ the background cell size and $c_d$ the dilatational wave speed. A
	safety factor $0<\mathrm{CFL}\le 1$ is normally applied,
	$\Dt=\mathrm{CFL}\cdot h/c_d$.

	\paragraph{Parabolic (diffusion) limit of the phase field.}
	Treating~\eqref{eq:evolution} as a diffusion equation with diffusivity
	$D_\phi=\tfrac{\ell\,G_c}{2c_w\,\eta}$, the forward--Euler update is stable only
	if
	\begin{equation}\label{eq:pf_stab}
		\Dt \;\le\; \Dt_{\mathrm{crit}}^{\,\phi} \;\sim\; \frac{h^2}{2\,n\,D_\phi}
		\;=\; \frac{c_w\,\eta\,h^2}{n\,\ell\,G_c},
	\end{equation}
	with $n$ the spatial dimension. Equation~\eqref{eq:pf_stab} shows the role of
	the viscosity $\eta$: a larger $\eta$ relaxes the phase--field time--step
	restriction (and brings the response closer to rate--independent fracture), but
	also slows the rate at which the crack can grow, so that $\eta$ must be small
	enough not to artificially retard the physical crack speed. In practice $\eta$
	is chosen so that $\Dt_{\mathrm{crit}}^{\,\phi}\gtrsim\Dt_{\mathrm{crit}}^{\,u}$,
	i.e.\ so that the elastodynamic CFL condition governs and the phase field can be
	advanced with the same $\Dt$ as the momentum update,
	\begin{equation}\label{eq:eta_choice}
		\eta \;\gtrsim\; \frac{n\,\ell\,G_c}{c_w\,h^2}\,\Dt_{\mathrm{crit}}^{\,u}
		\;=\; \frac{n\,\ell\,G_c}{c_w\,h\,c_d}.
	\end{equation}

	\paragraph{Consistency.}
	Because $\omega_{\mathrm{diss}}=\eta\dot\phi\to 0$ as $\dot\phi\to 0$, the
	stationary solutions of the explicit scheme coincide with the equilibrium
	solutions of the original quasi--static variational theory: the viscous term
	affects only the transient path towards equilibrium, not the converged
	crack pattern, provided $\eta$ is sufficiently small. This is the sense in which
	the dissipative regularisation is a purely \emph{algorithmic} device enabling
	the explicit time march while preserving the underlying fracture physics.

	% =====================================================================
	\section{Summary}

	Starting from the variational phase--field fracture theory, a single
	constitutive enrichment---the dissipative phase--field micro--force
	$\omega_{\mathrm{diss}}=\eta\,\dot\phi$---was introduced. This converts the
	elliptic damage stationarity condition into a parabolic Allen--Cahn evolution
	law~\eqref{eq:evolution}, which can be integrated by an explicit forward--Euler
	update on a lumped pseudo--mass~\eqref{eq:mpm_update}. The update was embedded in
	the explicit Update--Stress--Last MPM cycle, yielding the fully explicit
	\eng{USL--MPM--PF} solution procedure of Algorithm~\ref{alg:uslmpm_pf}, in which
	displacement and phase field are advanced within a single time loop subject to
	the combined CFL/diffusion time--step restriction. Both AT1 and AT2 crack
	density models are supported through the functions $w(\phi)$ and the
	normalisation $c_w$, and damage irreversibility is enforced through the
	monotone history field $\mathcal{H}$.


	% =====================================================================

	\begin{thebibliography}{9}

	\bibitem{nguyen2023mpm}
	V.P.\ Nguyen, A.\ de Vaucorbeil, S.\ Bordas,
	\textit{The Material Point Method: Theory, Implementations and Applications}.
	Springer, Scientific Computation, 2023.\\
	\href{https://doi.org/10.1007/978-3-031-24070-6}{https://doi.org/10.1007/978-3-031-24070-6}

	\bibitem{kristensen2021pff}
	P.K.\ Kristensen, C.F.\ Niordson, E.\ Mart\'inez-Pa\~neda,
	An assessment of phase field fracture: crack initiation and growth.
	\textit{Philosophical Transactions of the Royal Society A}, 379(2203):20210021, 2021.\\
	\href{https://doi.org/10.1098/rsta.2021.0021}{https://doi.org/10.1098/rsta.2021.0021}

	\bibitem{hu2023explicit}
	Z.\ Hu, Z.\ Zhang, X.\ Zhou, X.\ Cui, H.\ Ye,
	Explicit phase-field material point method with the convected particle domain interpolation for impact/contact fracture in elastoplastic geomaterials.
	\textit{Computer Methods in Applied Mechanics and Engineering}, 405:115851, 2023.\\
	\href{https://doi.org/10.1016/j.cma.2022.115851}{https://doi.org/10.1016/j.cma.2022.115851}

	\end{thebibliography}

\end{document}
